\subsection{Related Work}

A standard way to handle the composition $h\circ\calA$ in
Problem~\eqref{prob:main} is to introduce the auxiliary variable
$\by=\calA(\bx)$ and solve the equality-constrained problem
\begin{equation}\label{prob:split}
  \min_{\bx,\by}\; f(\bx)+h(\by)\quad\text{s.t.}\quad \calA(\bx)=\by.
\end{equation}
Most methods for~\eqref{prob:split} are double-loop: prox-linear
methods~\cite{Drusvyatskiy2019}, quadratic penalty
methods~\cite{kong2019complexity,lin2022complexity}, and inexact or proximal
augmented Lagrangian methods (ALMs)~\cite{sahin2019inexact,li2021rate,xie2021proximal}
solve a subproblem approximately at every outer iteration, and their
performance depends on the inner solver and its termination criterion.
Single-loop methods avoid this overhead by taking one explicit gradient or
proximal step per variable at each iteration. Without additional regularity
assumptions, however, existing single-loop methods for
Problem~\eqref{prob:main} attain $\mathcal O(\epsilon^{-3})$ complexity only by
scaling a parameter with $\epsilon$: Moreau-envelope
smoothing~\cite{bohm2021variable} takes the smoothing parameter of order
$\epsilon$, and single-loop penalty and linearized ALM/ADMM
schemes~\cite{hajinezhad2019perturbed,luSingleLoopGradientDescent2022,hallak2023adaptive,alacaoglu2024complexity,sun2024dual_descent_alm_admm,elbourkhissi2025perturbed,shi2025momentum}
take the penalty or dual perturbation parameter of order $\epsilon^{-1}$. In
both cases the smoothness constant of the primal problem grows as
$\epsilon^{-1}$, and the primal stepsize shrinks to $\Theta(\epsilon)$.
Fixed-penalty methods with $\mathcal O(\epsilon^{-2})$ complexity require a
polyhedral error bound for affine constraints~\cite{zhang2022global} or a
uniform LICQ for smooth nonlinear
constraints~\cite{elbourkhissi2025linearized,liu2026fixed}; we assume neither
condition for Problem~\eqref{prob:main}.

Eliminating $\by$ from the Lagrangian of~\eqref{prob:split} gives the
nonconvex--linear minimax problem
\begin{equation}\label{eq:ncl}
  \min_{\bx\in\RR^n}\max_{\bq\in\RR^m}\;
  \calL(\bx,\bq) = \Phi(\bx,\bq)-h^*(\bq), \qquad
  \Phi(\bx,\bq):=f(\bx)+\langle\bq,\calA(\bx)\rangle,
\end{equation}
whose game-stationary (GS) points coincide with the KKT points
of~\eqref{prob:split}. For a fixed penalty $\rho>0$, eliminating $\by$ from the
augmented Lagrangian instead gives a smooth nonconvex--concave (NC--C) minimax
problem, and a single-loop linearized ALM is a gradient descent--ascent (GDA)
method on it (Section~\ref{subsec:reduced-al}). For smooth NC--C minimax
problems, two-timescale GDA requires $\mathcal O(\epsilon^{-6})$
iterations~\cite{lin2020gradient}. Primal smoothing (Smoothed
GDA~\cite{zhang2020smoothed}) and dual perturbation (Perturbed
GDA~\cite{lu2020hybrid,xu2023unified}) each reduce this to
$\mathcal O(\epsilon^{-4})$, and Perturbed Smoothed GDA, which combines the two,
attains $\mathcal O(\epsilon^{-3})$ with a dual regularization parameter of
order $\epsilon^{-1}$~\cite{li2026smoothing}. Li et al.~\cite{li2026smoothing}
also prove matching lower bounds of $\Omega(\epsilon^{-4})$ and
$\Omega(\epsilon^{-3})$ for these two algorithms. Single-loop methods designed
for the nonconvex--linear structure likewise attain
$\mathcal O(\epsilon^{-3})$ with an $\epsilon$-dependent dual
regularization~\cite{pan2021efficient,he2024approximation,xu2026rada}.

\paragraph{Our contributions}
In every $\mathcal O(\epsilon^{-3})$ guarantee above, either the penalty
parameter or the dual regularization parameter is of order $\epsilon^{-1}$.
For ALMs, this growth matters: a fixed penalty is the main advantage of ALMs
over penalty methods, because it keeps the primal smoothness constant bounded
as $\epsilon\to0$. Existing lower bounds do not settle whether the growth is
necessary. The bounds in~\cite{mahdavinia2022tight,li2026smoothing} hold for
individual algorithms under prescribed parameters, and the problem-class
bounds for smooth nonconvex or affinely constrained
problems~\cite{carmon2020lower,liu2026lower} do not address the choice of the
penalty. This raises the following question:
\begin{quote}
  \textbf{(Q1)}
  \emph{Without additional regularity assumptions, must a single-loop
  linearized ALM let some parameter grow with $1/\epsilon$ to attain
  $\mathcal O(\epsilon^{-3})$ complexity? If so, must the penalty parameter
  carry this growth?}
\end{quote}

In this paper, we provide a precise characterization addressing
\textbf{(Q1)}. We introduce LRALM, a family $\calC_\theta$ of single-loop
linearized ALMs with a penalty $\rho$, a dual regularization parameter $s$, and
primal and dual centers updated with constant center stepsizes
(Section~\ref{sec:algorithm-framework}). The family contains Smoothed GDA and
Perturbed Smoothed GDA, and every member has the effective penalty $\rho+s$
(Lemma~\ref{lem:effective-penalty}). We first show that if
$\rho+s\le c_1/\epsilon$, every member of $\calC_\theta$ requires
$\Omega(\epsilon^{-4})$ iterations to find an $\epsilon$-KKT point
(Theorem~\ref{thm:lower-bound}), and that a member of $\calC_\theta$ attains
this bound (Theorem~\ref{thm:upper-beta1}). We further show that a fixed $\rho$
with $s=\Theta(\epsilon^{-1})$ attains $\mathcal O(\epsilon^{-3})$
(Theorem~\ref{thm:upper-beta0}), and that every member of $\calC_\theta$ with
these parameters requires $\Omega(\epsilon^{-3})$ iterations
(Theorem~\ref{thm:lower-bound-fixed-rho}). Hence the answer to the first part
of \textbf{(Q1)} is positive and the answer to the second part is negative:
some parameter must grow as $1/\epsilon$, and the dual regularization
parameter can carry this growth while $\rho$ stays fixed.

